\weeknum{10}
\topic[Finite Differences]{Geothermics: Finite Differences}

\workshopstart

\vspace{-0.75em}
\section*{Preparation}

\begin{description}
  \item[Pre-reading:] Course Notes, Ch. 18
  \item[Lecture videos:]
  \item[Notes:]
\end{description}

\section*{Purpose}

This workshop builds intuition for finite difference methods through physical reasoning and interactive activities. Students first derive the heat equation from energy conservation, then experience how numerical models work by acting as nodes that update temperature based on their neighbors.  The workshop then explores how boundary conditions represent geological assumptions and how numerical stability reflects limits on how heat can physically propagate.  By the end, students should understand that finite difference methods are simply local energy balance, that heat flow emerges from neighbor interactions, and that both boundary conditions and stability have direct physical meaning.

\section*{Learning Outcomes}
Students will:
\begin{itemize}
    \item Understand finite difference equations as local statements of energy conservation.
    \item Develop intuition for how numerical solutions evolve through neighbor interactions.
    \item Interpret boundary conditions as geological assumptions.
    \item Understand numerical stability as a limit on physical information transfer.
\end{itemize}

\section*{Session Flow (\timelinelength\ min)}

\timeline[slot]{Warm-up --- What breaks down?}{10}

Think back over the analytic solutions in Weeks 8 and 9 (steady-state geotherm, half-space cooling, plate cooling). Discuss a geological situation where each analytic solution would fail. What assumption(s) breaks down?

\timeline[slot]{Mini-lecture --- From energy conservation to finite differences}{15}

\begin{itemize}
    \item The heat equation as a statement of local energy balance on a small rock volume.
    \item Discretization: replacing continuous derivatives with differences between neighboring nodes.
    \item The finite difference stencil: each node updates based on its neighbors.
    \item Stability: why the time step cannot be arbitrarily large ($\Delta t < \Delta x^2 / 2\kappa$).
\end{itemize}

What makes geological systems difficult to solve analytically? Consider:
\begin{itemize}
    \item Analytic solutions require simplifications:
    \begin{itemize}
        \item constant properties
        \item simple geometries
        \item simple boundary conditions
    \end{itemize}
    \item Real geology violates these assumptions.
\end{itemize}

\timeline[item]{From Conservation of Energy to Finite Differences}{30}

Finite difference equations arise from applying conservation of energy to small
rock volumes.

\begin{itemize}
    \item Nodes represent local rock volumes.
    \item Heat flux enters and leaves each node.
    \item Heat production adds energy locally.
\end{itemize}

The equation is \emph{not a numerical trick}.  The Laplacian as a measure of curvature (difference from neighbors).  This is just bookkeeping of energy.

\timeline[item]{Human Finite Difference Solver}{30}

Numerical models evolve through local interactions between neighboring nodes. We'll do an exercise with you as the computer to see how finite difference solutions work.

We'll explore
\begin{itemize}
    \item smoothing of temperature
    \item spreading of heat
    \item role of gradients
\end{itemize}

\timeline[slot]{Break}{10}

\timeline[item]{Boundary Conditions as Geological Assumptions}{30}

Boundary conditions determine what physical system the model represents. There are two fundamental types:
\begin{itemize}
    \item Fixed value (Dirichlet)
    \item Fixed flux (Neumann)
\end{itemize}
We'll also look at two additional types (Robin and Cauchy) that are combinations or independent applications of these types.

\timeline[item]{Stability and Information Transfer}{25}

Stability limits impose a ``speed limit'' on thermal information:
\[
    \Delta t < \frac{\Delta x^2}{2\kappa}
\]
Heat propagates gradually through diffusion---numerical schemes must respect this physical constraint. We'll explore:
\begin{itemize}
    \item small vs large time steps
    \item what happens when updates are too large
    \item stability = causality
    \item information cannot skip nodes
\end{itemize}

Instability is not just numerical---it is physically impossible behavior.

\timeline[slot]{Wrap-up}{10}

\textbf{Key connections}
\begin{itemize}
    \item Finite differences = local energy balance + neighbor interaction
    \item Boundary conditions = geological assumptions
    \item Stability = physical constraint on information transfer
\end{itemize}

\textbf{Discussion}
\begin{itemize}
    \item How does the numerical method compare to analytic solutions?
    \item What advantages do numerical models provide?
    \item Where might these methods be applied beyond heat flow?
\end{itemize}

\timeline[slot]{Course Conclusion}{20}

This is the final workshop of the course. Rather than a preview, use this time to connect ideas across all ten weeks.

\begin{description}
    \item [Fields and gradients] (Week 1): all geophysical observations are measurements of fields; gradients tell us where physical property contrasts occur.
    \item [Physical properties] (Week 2): the link between measurement and geology lives in physical properties and their mixing laws.
    \item [Inversion] (Week 3): going from data to model requires confronting non-linearity, non-uniqueness, and the role of prior information.
    \item [Gravity and isostasy] (Weeks 4--5): density contrasts at depth shape both the gravity field and the long-term mechanical balance of the lithosphere.
    \item [Magnetics and Fourier methods] (Weeks 6--7): vector fields, polarity, depth--wavelength relationships, and spectral filtering.
    \item [Geothermics] (Weeks 8--9): heat flow connects surface measurements to deep thermal structure via conduction, diffusion, and phase changes.
    \item [Finite differences] (Week 10): numerical methods let us move beyond analytic limits to model the real, heterogeneous Earth.
\end{description}

\workshopend
\notepage
\notepage

\subimport{activities/}{activity_finite_differences}
\subimport{activities/}{activity_human_computer}
\subimport{activities/}{activity_geologic_boundaries}
\subimport{activities/}{activity_stability}